% ============================================================
% related_work.tex — Sección II: Trabajos Relacionados
% Estado: Desarrollado completamente
% ============================================================

\section{Trabajos Relacionados}
\label{sec:related_work}

Esta sección analiza las líneas de investigación relacionadas con la propuesta, organizadas en cuatro ejes temáticos fundamentales.

\subsection{Retrieval-Augmented Generation (RAG)}
\label{subsec:rw_rag}

El paradigma \rag{} fue formalizado por Lewis et al.~\cite{lewis2020retrieval}, quienes propusieron la combinación de un retriever con un generator para tareas de NLP intensivas en conocimiento, demostrando mejoras en factualidad y relevancia frente a modelos puramente paramétricos. Guu et al.~\cite{guu2020realm} introdujeron REALM, integrando la recuperación directamente en el pre-entrenamiento del modelo de lenguaje. Borgeaud et al.~\cite{borgeaud2022improving} escalaron la idea mediante RETRO, demostrando que modelos con acceso a bases de conocimiento externas de trillones de tokens alcanzan rendimiento competitivo con modelos significativamente más grandes.

Izacard y Grave~\cite{izacard2021leveraging} propusieron Fusion-in-Decoder, concatenando múltiples pasajes recuperados en el decoder. Asai et al.~\cite{asai2024selfrag} introdujeron Self-RAG, un framework que enseña al modelo a recuperar, generar y criticar sus propias respuestas mediante tokens de reflexión. Khattab y Zaharia~\cite{khattab2020colbert} desarrollaron ColBERT, un modelo de interacción tardía para búsquedas eficientes mediante representaciones contextualizadas. Gao et al.~\cite{gao2024retrieval} proporcionaron un survey comprehensivo del paradigma \rag{}, categorizando las aproximaciones según sus componentes de recuperación, aumentación y generación.

A pesar de estos avances, la aplicación de \rag{} en dominios biomédicos especializados permanece subexplorada. La mayoría de los sistemas operan sobre corpora genéricos y no incorporan mecanismos de adquisición continua ni memoria persistente.

\subsection{NLP biomédico y recuperación semántica}
\label{subsec:rw_bionlp}

Lee et al.~\cite{lee2020biobert} desarrollaron \biobert{}, un modelo basado en BERT pre-entrenado con artículos de PubMed y PMC, demostrando mejoras sustanciales en tareas de extracción biomédica y question answering. Gu et al.~\cite{gu2021domain} propusieron PubMedBERT, argumentando que el pre-entrenamiento desde cero con vocabulario especializado supera al pre-entrenamiento continuo. Neumann et al.~\cite{neumann2019scispacy} presentaron \scispacy{}, una extensión de spaCy optimizada para texto biomédico con modelos para NER y vinculación a ontologías como UMLS.

Luo et al.~\cite{luo2022biogpt} introdujeron BioGPT, un modelo generativo pre-entrenado para el dominio biomédico. Jin et al.~\cite{jin2019pubmedqa} crearon PubMedQA, un benchmark de referencia para question answering biomédico. Xiong et al.~\cite{xiong2021approximate} propusieron ANCE para mejorar la calidad de recuperación densa en dominios especializados. Nentidis et al.~\cite{nentidis2023overview} reportaron los avances del desafío BioASQ, evidenciando los retos persistentes en QA biomédico.

\subsection{Bases de datos vectoriales y memoria persistente}
\label{subsec:rw_vectordb}

Johnson et al.~\cite{johnson2021billion} desarrollaron \faiss{}, estableciendo los fundamentos algorítmicos (IVF, PQ, HNSW) para la búsqueda eficiente de vecinos más cercanos a escala. Reimers y Gurevych~\cite{reimers2019sentence} introdujeron Sentence-BERT, ampliamente adoptado para la representación vectorial de documentos. Sistemas como \chromadb{}, Milvus y Weaviate han democratizado el acceso a bases de datos vectoriales con persistencia nativa.

Park et al.~\cite{park2023generative} exploraron la integración de memoria a largo plazo en agentes generativos, donde las representaciones vectoriales de interacciones pasadas se almacenan para mantener coherencia contextual entre sesiones, una capacidad fundamental para la arquitectura propuesta en este trabajo.

\subsection{Sistemas para etnobotánica y farmacognosia}
\label{subsec:rw_ethnobotany}

Bussmann y Sharon~\cite{bussmann2006traditional} documentaron la flora medicinal del norte de Perú, proporcionando una base de datos de referencia. Heinrich et al.~\cite{heinrich2020ethnopharmacology} revisaron las aproximaciones computacionales a la etnofarmacología, identificando oportunidades para la aplicación de técnicas de IA. Gonzales~\cite{gonzales2012ethnobiology} proporcionó una revisión comprehensiva de la evidencia sobre \textit{Lepidium meyenii}. Chandak et al.~\cite{chandak2023building} exploraron grafos de conocimiento biomédico para la representación de relaciones entre compuestos y actividades biológicas.

\subsection{Posicionamiento de la propuesta}
\label{subsec:rw_positioning}

La Tabla~\ref{tab:related_work_comparison} resume el posicionamiento de \sirca{} frente a los sistemas revisados. Ningún sistema existente integra simultáneamente las cinco capacidades propuestas: pipeline \rag{} biomédico, memoria vectorial persistente, retrieval híbrido, adquisición continua y grounding documental, específicamente orientado al dominio etnobotánico.

% ============================================================
% related_work_table.tex — Tabla comparativa de trabajos relacionados
% ============================================================

\begin{table*}[!t]
\centering
\caption{Análisis comparativo de sistemas relacionados frente a las capacidades de \sirca{}.}
\label{tab:related_work_comparison}
\renewcommand{\arraystretch}{1.2}
\footnotesize
\begin{tabularx}{\textwidth}{l c c c c c c}
\toprule
\textbf{Sistema / Enfoque} & \textbf{RAG} & \textbf{Memoria} & \textbf{Retrieval} & \textbf{Adquisición} & \textbf{Grounding} & \textbf{Dominio} \\
 & \textbf{Biomédico} & \textbf{Persistente} & \textbf{Híbrido} & \textbf{Continua} & \textbf{Documental} & \textbf{Específico} \\
\midrule
RAG~\cite{lewis2020retrieval}                   & --          & --          & --          & --          & Parcial     & General \\
REALM~\cite{guu2020realm}                       & --          & --          & Denso       & --          & --          & General \\
Self-RAG~\cite{asai2024selfrag}                 & --          & --          & Denso       & --          & \checkmark  & General \\
ColBERT~\cite{khattab2020colbert}               & --          & --          & Denso       & --          & --          & General \\
BioBERT~\cite{lee2020biobert}                   & --          & --          & --          & --          & --          & Biomédico \\
BioGPT~\cite{luo2022biogpt}                     & --          & --          & --          & --          & --          & Biomédico \\
PubMedQA~\cite{jin2019pubmedqa}                 & --          & --          & --          & --          & Parcial     & Biomédico \\
S2ORC~\cite{lo2020s2orc}                        & --          & --          & --          & \checkmark  & --          & General \\
\midrule
\textbf{\sirca{} (propuesto)} & \checkmark  & \checkmark  & \checkmark  & \checkmark  & \checkmark  & \textbf{Etnobotánico} \\
\bottomrule
\end{tabularx}
\end{table*}

